\section{5. Conclusion}

\section{5. Conclusion}

In this paper, we propose a Unified Dual-Mask Physical Model for non-Lambertian light field (LF) depth estimation to address the fundamental failure of conventional Epipolar Plane Image (EPI) based methods under complex reflectance conditions. By abandoning the strict photo-consistency assumption, this work constructs a physically interpretable framework that explicitly models the intricate light transport mechanisms of diverse materials. Specifically, we introduce a novel paradigm that effectively integrates material reflection analysis with dual-mask physical modeling, enabling robust depth recovery across Lambertian, non-Lambertian, and mixed-material scenes. The core contributions of this work are summarized as follows:

\textit{ \textbf{MRI-Inspired Angular Frequency Analysis for Material Parsing:} We developed an angular frequency analysis technique inspired by MRI to parse material reflection types. This approach overcomes the inherent bottleneck of distinguishing complex reflectance properties from limited angular samples, providing a physically meaningful and computationally efficient material classification prior that lays a reliable foundation for subsequent depth estimation in non-Lambertian scenes.
} \textbf{Dual-Mask Physical Modeling:} We formulated a dual-mask physical modeling strategy encompassing both medium and angular dimensions. By unifying the modeling of Lambertian, non-Lambertian, and mixed scenes from a physical optics perspective, this mechanism resolves the lack of physical interpretability in conventional data-driven models, empowering the network to adaptively process mixed materials and complex illumination conditions.
\textbf{Component-Aware Adaptive Depth Fusion:} We designed a component-aware adaptive depth estimation fusion strategy. This mechanism substantially resolves the fatal flaw of traditional EPI methods—where depth prediction degenerates into mere texture copying in non-Lambertian regions due to the violation of the photo-consistency assumption—thereby significantly enhancing depth estimation accuracy and robustness in mixed and non-Lambertian scenarios.

A series of experiments conducted on diverse light field datasets validate the superiority of our Unified Dual-Mask Physical Model across multiple dimensions, demonstrating state-of-the-art performance particularly in challenging non-Lambertian and mixed-material regions. Spatially-varying BRDFs and complex specularities can also be effectively handled by our method, while this is largely unattainable using conventional EPI-based or purely data-driven approaches. Looking forward, the proposed physical priors and dual-mask mechanism hold substantial potential for broader 3D vision applications, such as novel view synthesis and dynamic light field reconstruction, enabling for more physically grounded and generalizable computational imaging systems.
